% Quadrature: the waveform, the four states, where the counts come from, and
% what the index adds. Laid out in bands with nothing beside anything, because
% this figure carries three separate ideas and they crowd each other easily.
\begin{tikzpicture}[font=\sffamily\scriptsize, x=1cm, y=1cm]

% ================================================================ band 1: signal
\node[signame] at (-0.3,2.5) {A};
\node[signame] at (-0.3,1.4) {B};
\node[signame] at (-0.3,0.3) {Z};

\draw[signal] (0,2.5) -- (1,2.5) -- (1,2.1) -- (2,2.1) -- (2,2.5) -- (3,2.5)
              -- (3,2.1) -- (4,2.1) -- (4,2.5) -- (5,2.5) -- (5,2.1) -- (6,2.1)
              -- (6,2.5) -- (7,2.5) -- (7,2.1) -- (8,2.1);
\draw[signalb] (0,1.4) -- (0.5,1.4) -- (0.5,1.0) -- (1.5,1.0) -- (1.5,1.4)
               -- (2.5,1.4) -- (2.5,1.0) -- (3.5,1.0) -- (3.5,1.4) -- (4.5,1.4)
               -- (4.5,1.0) -- (5.5,1.0) -- (5.5,1.4) -- (6.5,1.4) -- (6.5,1.0)
               -- (7.5,1.0) -- (7.5,1.4) -- (8,1.4);
\draw[signal] (0,0.3) -- (4,0.3) -- (4,0.7) -- (4.25,0.7) -- (4.25,0.3) -- (8,0.3);
\node[tlabel, anchor=west] at (4.45,0.6) {one index pulse every N counts};

\foreach \x in {0.5,1,1.5,2,2.5,3,3.5,4,4.5,5,5.5,6,6.5,7,7.5}{
  \draw[tick] (\x,0.2) -- (\x,2.8);
  \node[sample] at (\x,2.95) {};
}
\node[tlabel, anchor=west] at (8.2,2.95) {every marked edge is one count};
\node[tlabel, anchor=west] at (8.2,2.35) {four per cycle: the mode this node uses};

% ================================================================ band 2: states
\node[chlabel, anchor=east] at (-0.3,-0.4) {state, A then B};
\foreach \x/\s in {0.25/10, 0.75/11, 1.25/01, 1.75/00, 2.25/10, 2.75/11,
                   3.25/01, 3.75/00, 4.25/10, 4.75/11, 5.25/01, 5.75/00,
                   6.25/10, 6.75/11, 7.25/01, 7.75/00}{
  \node[font=\sffamily\tiny\ttfamily] at (\x,-0.4) {\s};
}
\draw[-Latex, draw=linecol] (0.25,-1.0) -- (3.75,-1.0);
\node[tlabel, anchor=west] at (3.95,-1.0) {forward: 10, 11, 01, 00, and repeat};
\draw[-Latex, draw=red!70!black] (7.75,-1.6) -- (4.25,-1.6);
\node[tlabel, anchor=east, text=red!60!black] at (4.05,-1.6) {reverse: the same four states, the other way round};

% ================================================================ band 3: modes
\node[chlabel, anchor=west] at (0,-2.4) {the three modes this part offers, and only these three};
\foreach \i/\m/\d/\fillc in {0/{count on A's edges}/{two counts per cycle}/white,
                             1/{count on B's edges}/{two counts per cycle}/white,
                             2/{count on both}/{four counts per cycle, and free}/hwcol}{
  \node[regcell, minimum width=34mm, text width=32mm, anchor=south west,
        fill=\fillc] at (\i*3.6,-3.3) {\m\\{\tiny \d}};
}

% ================================================================ band 4: the reading
\node[chlabel, anchor=north west, text width=54mm] at (0,-3.9)
  {The mapping from those three names to two and four counts per cycle is
   standard reference-manual semantics. It was not read in the manual for this
   part during the research for this volume, so the chapter confirms it on the
   bench by generating a known number of cycles and counting what arrives, which
   is better evidence than a citation anyway.};

\node[umlnote, text width=54mm, anchor=north west] at (5.8,-3.9)
  {\textbf{What the index is for.} A quadrature decoder cannot detect an edge it
   never saw: a dropped edge is a permanent one-count error and nothing in the
   signal reveals it. The index pulse is the only thing that turns a relative
   count back into an absolute position, which is why an incremental encoder
   without an index is a poor choice for a joint, and why this part having no
   hardware index support is worth a whole step of the chapter.};
\end{tikzpicture}
